\chapter{Discussione}
\label{chap:discussione}

\section{Confronto con il paper di riferimento}

Il paper BoNet~\cite{escobar2019bonet} riporta, per la keypoint
estimation sul dataset RHPE, un mAP@[.5:.95] di 0.814 con bounding box
ground truth e di 0.808 con bounding box predetta. La
Tabella~\ref{tab:confronto_bonet} confronta questi valori con i
risultati ottenuti in questo lavoro, sia sul sottoinsieme da 250
immagini sia con \texttt{gt\_crop} riallenato sul dataset completo con
data augmentation (Sez.~\ref{sec:retraining_gt_crop_full}).

\begin{table}[H]
\centering
\caption{Confronto tra i valori di BoNet e i risultati di questo lavoro
(modello sempre \texttt{gt\_crop}, Sez.~\ref{sec:retraining_gt_crop_full});
\texttt{gt\_crop}/\texttt{det\_crop} in tabella indicano qui la
modalità di valutazione.}
\label{tab:confronto_bonet}
\smallskip
\small
\begin{tabular}{llcc}
\toprule
Fonte & Condizione & mAP@[.5:.95] & Divario \\
\midrule
BoNet & bbox ground truth & 0.814 & \\
\texttt{gt\_crop} & sottoinsieme & 0.7037 & 0.110 \\
\texttt{gt\_crop} & dataset completo + aug & 0.7699 & 0.044 \\
\midrule
BoNet & bbox predetta & 0.808 & \\
\texttt{det\_crop} & sottoinsieme & 0.7205 & 0.088 \\
\texttt{det\_crop} & dataset completo + aug & 0.7913 & 0.017 \\
\bottomrule
\end{tabular}
\end{table}

Il valore in modalità \texttt{det\_crop} supera quello in modalità
\texttt{gt\_crop} per un motivo già discusso in
Sez.~\ref{sec:retraining_gt_crop_full} (l'augmentation avvantaggia in
modo particolare gli input con un margine maggiore attorno alla mano),
non per un effetto isolato o casuale. Interpretare questi divari
richiede comunque di considerare più fattori, discussi di seguito.

\subsection{Affidabilità del valore di riferimento}
\label{sec:affidabilita_bonet}

Un'analisi diretta della Tabella 2 del paper BoNet solleva un dubbio
sull'affidabilità del valore 0.814 come target specifico per RHPE. La
tabella riporta i risultati di hand pose estimation separatamente per
RSNA e RHPE, in tre condizioni (immagine intera, crop con bounding box
ground truth, crop con bounding box predetta). Mentre la riga
``immagine intera'' differisce sensibilmente tra i due dataset (73.0\%
RSNA contro 53.1\% RHPE, coerentemente con il fatto che le immagini
RHPE, a differenza di quelle RSNA, includono entrambe le mani e altri
elementi come etichette identificative), le righe ``bounding box ground
truth'' e ``bounding box predetta'' riportano valori identici, cifra
per cifra, tra i due dataset su tutte e quattro le metriche. Questa
coincidenza esatta su due righe, a fronte di una terza che differisce in
modo plausibile, è compatibile con un errore di copia nella tabella del
paper originale, più che con un risultato genuino. Non è stato possibile
risolvere il dubbio: il materiale supplementare del paper non riporta
una tabella alternativa, il repository ufficiale del
codice\footnote{\url{https://github.com/BCV-Uniandes/Bonet}} include
solo l'implementazione del terzo stadio (bone age assessment), non
della hand detection o della hand pose estimation, e nessun lavoro
successivo che cita questo valore lo riporta in modo indipendente. Il
valore 0.814 va quindi trattato con cautela, non come un target
RHPE-specifico verificato con certezza.

\subsection{Differenze architetturali e di training}
\label{sec:differenze_architetturali_bonet}

Anche assumendo che 0.814 sia un valore affidabile, il confronto resta
soggetto ad almeno due differenze non isolate in questo lavoro.

La prima è architetturale: BoNet utilizza Simple Baselines for Human
Pose Estimation~\cite{xiao2018simplebaselines}, mentre questo lavoro
utilizza Keypoint R-CNN~\cite{he2017maskrcnn}. Un divario di prestazioni
tra le due architetture è documentato indipendentemente sul benchmark
standard della pose estimation, il task di keypoint detection del
dataset COCO~\cite{lin2014coco} (un compito diverso, la posa umana su
fotografie anziché mani radiografiche, ma quello su cui entrambi i paper
riportano i propri risultati): Simple Baselines con backbone ResNet-50
raggiunge un AP di 70.4 su COCO val2017~\cite{xiao2018simplebaselines},
mentre Mask R-CNN con lo stesso backbone (ResNet-50-FPN), da cui deriva
Keypoint R-CNN, raggiunge un AP di 62.7 su COCO
test-dev~\cite{he2017maskrcnn} (le due valutazioni usano split diversi
del dataset, quindi il confronto numerico esatto va preso con cautela,
ma la direzione e l'ordine di grandezza del divario, 7-8 punti, sono
coerenti con quanto osservato in questo lavoro). La spiegazione
tipicamente proposta è che Simple Baselines opera su un crop già
ritagliato ad alta risoluzione con un decoder a deconvoluzione dedicato,
mentre Keypoint R-CNN predice i keypoint da una feature map ottenuta
tramite RoIAlign, storicamente meno precisa per la localizzazione fine.

La seconda è una discrepanza nella ricetta di training, mai isolata
prima d'ora. Il paper BoNet indica, per la keypoint estimation, solo il
numero di epoche (20) e l'uso dei ``parametri suggeriti'' da Simple
Baselines: ottimizzatore Adam, learning rate $10^{-3}$, batch size 128.
Poiché Keypoint R-CNN appartiene alla famiglia Faster R-CNN,
tradizionalmente addestrata con SGD, questo lavoro ha ereditato da
BoNet solo il numero di epoche, non il resto della ricetta:
l'addestramento usa SGD con batch size 4, un numero di aggiornamenti dei
pesi per epoca circa 32 volte più alto a parità di epoche nominali. ``20
epoche'' non rappresenta quindi lo stesso volume di addestramento nei
due casi.

L'introduzione della data augmentation descritta in
Sez.~\ref{sec:keypoint_estimation}, motivata proprio da questo
confronto, ha permesso di ridurre il divario da 0.063 a 0.044 (modalità
\texttt{gt\_crop}) senza modificare l'architettura né la ricetta di
ottimizzazione. Questo è compatibile con l'idea che il divario residuo
derivi da una combinazione delle due differenze appena descritte,
sebbene isolarne il contributo relativo richiederebbe esperimenti
dedicati, che restano fuori dallo scope di questo lavoro.

\section{Osservazioni metodologiche}

Il pattern osservato in Sez.~\ref{sec:risultati_keypoint}, un crollo
quasi totale delle prestazioni quando un modello viene valutato fuori
dal dominio visto in training, non è un artefatto specifico di questo
lavoro, ma un'istanza di un fenomeno più generale: un modello di
apprendimento supervisionato non impara un compito in astratto, impara
a risolverlo assumendo implicitamente che i dati futuri seguano la
stessa distribuzione di quelli visti in addestramento
(Sez.~\ref{sec:classificazione}). Più dati e data augmentation
attenuano parzialmente questo crollo, senza eliminarlo: lo stesso
\texttt{gt\_crop} valutato in modalità \texttt{full}, riallenato sul
dataset completo con augmentation, passa da 0.086 a 0.40
(Sez.~\ref{sec:retraining_gt_crop_full}), un miglioramento reale ma
ancora lontano dal picco del modello nel proprio dominio (0.77).

Nel caso di \texttt{gt\_crop} e
\texttt{mirrored}, questa assunzione implicita è particolarmente
rigida: il target di bounding box è sempre l'intero crop, quindi il
modello non solo non ha mai visto immagini con sfondo, ma non ha mai
avuto un incentivo ad apprendere a distinguere la mano dal contesto
circostante.

Un confronto più specifico è compatibile con questa lettura.
\texttt{full} e \texttt{mirrored\_full} sono gli unici due modelli
addestrati sull'immagine intera; quando entrambi vengono valutati fuori
da questo dominio, su crop, \texttt{mirrored\_full} supera \texttt{full}
in modo pressoché sistematico su quasi tutte le combinazioni di split e
metrica (ad esempio, sul test set, mAP@[.5:.95] in modalità
\texttt{gt\_crop}: 0.527 contro 0.515; in modalità \texttt{det\_crop}:
0.561 contro 0.541). I risultati sono compatibili con l'ipotesi che il
raddoppio dei dati di training introdotto dallo specchiamento, pensato
per predire keypoint su entrambe le mani, abbia anche un effetto
positivo secondario sulla capacità di generalizzazione del modello; non
è però possibile, con gli esperimenti disponibili, isolare questo
effetto da spiegazioni alternative (ad esempio un semplice aumento
della quantità di dati, indipendente dallo specchiamento in sé), poiché
non è stato addestrato un modello di controllo su 500 immagini non
specchiate.

Questo suggerisce una lettura pratica oltre a quella puramente
numerica: sebbene Keypoint R-CNN sia di per sé un'architettura
integrata capace di detection e keypoint estimation insieme, la testa
di detection di \texttt{gt\_crop} e \texttt{mirrored} non ha mai
imparato a localizzare la mano in un contesto ampio, essendo stata
addestrata con un target di bounding box sempre coincidente con
l'intero crop. In pratica, questi due modelli restano quindi
utilizzabili solo se a monte è disponibile una bounding box prodotta
altrove (dal ground truth o da un sistema di detection dedicato, come
nella modalità \texttt{det\_crop}): non è un limite architetturale di
Keypoint R-CNN, ma una conseguenza diretta di come sono stati
addestrati. \texttt{mirrored\_full}, addestrato invece sull'intera
immagine, non ha questo vincolo, ed è l'unica variante in grado di
funzionare in autonomia end-to-end: nessuna delle varianti allenate su
crop lo supera in modo netto nel confronto di picco
(\texttt{gt\_crop}: 0.704, sostanzialmente alla pari;
\texttt{mirrored}: 0.640, l'unica a restare visibilmente indietro,
Sez.~\ref{sec:risultati_keypoint}), quindi il vantaggio in termini di
robustezza operativa non comporta un sacrificio significativo in
termini di prestazioni.

\section{Limiti del lavoro}

Un primo limite riguarda la mano destra: le quattro strategie
confrontate in Sez.~\ref{sec:esplorazione_dx} sono state valutate solo
tramite ispezione visiva su un campione ridotto di immagini, poiché il
dataset RHPE annota esclusivamente la mano sinistra. Non è quindi
possibile fornire una misura quantitativa dell'accuratezza su questa
mano, né confrontarla direttamente con le prestazioni ottenute sulla
sinistra.

In secondo luogo, l'intera estensione alla mano destra si fonda
sull'assunzione che le due mani siano approssimabili tramite uno
specchiamento orizzontale (Sez.~\ref{sec:mirroring}). Si tratta di una
semplificazione geometrica: non tiene conto di eventuali asimmetrie
reali tra le due mani, né di differenze di rotazione o posizionamento
rispetto all'asse di specchiamento, e la sua validità non è stata
verificata clinicamente o anatomicamente.

Infine, tutti i modelli sono stati addestrati e valutati su un'unica
fonte di dati, il dataset RHPE, raccolto da una singola coorte con un
protocollo di acquisizione specifico. Si tratta, nella sostanza, dello
stesso limite discusso per BoneXpert in
Sez.~\ref{sec:approcci_automatici}: cambiare la popolazione di
riferimento, da danese a colombiana, non elimina il problema
strutturale di validare un modello su un'unica coorte. Il comportamento
dei modelli proposti in questo lavoro su dati provenienti da
popolazioni, protocolli di acquisizione o apparecchiature diverse resta
quindi non verificato.
